% Solution: Munoz-Garcia & Toro-Gonzalez (2019) Ch. 2 Exercise 7 (Cournot duopoly). Input by \solitem.
% The book prints the answer (pp. 48-50); this one is cross-checked against it. Check: check_L2.py.
\begin{exsol}{\MG{} Ch.~2 Exercise 7}{48}{Foundations}{2}
\asked Find the Nash equilibrium of the Cournot duopoly with inverse demand $p(Q) = a - bQ$,
$a > c$, $b > 0$, and common constant marginal cost $c$.

\method Cournot's model as a simultaneous game in quantities (Lecture~2, slides~20--25;
Tadelis \S5.2.3, pp.~87--89). Firm~$i$'s profit is concave in its own quantity, so its best
response solves the first-order condition; the Nash equilibrium is where the two best responses
cross (Lecture~2, slide~23).

\solution
\textbf{Game.} Players $1, 2$; strategies $q_i \ge 0$; payoff
\[
  \pi_i(q_i, q_j) = \bigl(a - b(q_i + q_j)\bigr) q_i - c\, q_i .
\]
\textbf{Best response.} Take $q_j$ as given.
\[
  \frac{\partial \pi_i}{\partial q_i} = a - 2b q_i - b q_j - c = 0
  \quad\Longrightarrow\quad
  q_i = \frac{a - c}{2b} - \frac12\, q_j .
\]
The second derivative is $-2b < 0$, so this is a maximum. The expression is negative when
$q_j > (a - c)/b$; since $q_i \ge 0$,
\[
  BR_i(q_j) = \max\Bigl\{0,\ \frac{a - c}{2b} - \frac{q_j}{2}\Bigr\},
\]
a line that starts at the monopoly output $(a - c)/(2b)$ when $q_j = 0$, falls with slope $-\tfrac12$
(quantities are strategic substitutes) and hits zero at $q_j = (a-c)/b$. The same holds for firm~$j$.

\textbf{Crossing.} Write $k = (a - c)/(2b)$, so $q_1 = k - q_2/2$ and $q_2 = k - q_1/2$ on the
interior branch. Substitute the second into the first:
\[
  q_1 = k - \tfrac12\bigl(k - \tfrac12 q_1\bigr) = \tfrac k2 + \tfrac{q_1}{4}
  \;\Longrightarrow\; \tfrac34\, q_1 = \tfrac k2
  \;\Longrightarrow\; q_1 = \tfrac{2k}{3} = \frac{a - c}{3b},
\]
and then $q_2 = k - k/3 = 2k/3$ as well. Both quantities are positive (since $a > c$), so the
interior branch of $BR_i$ is the right one at the solution. Aggregate output, price and profit follow:
\[
  Q^* = \frac{2(a - c)}{3b}, \qquad
  p^* = a - bQ^* = \frac{a + 2c}{3}, \qquad
  \pi^* = (p^* - c)\,q^* = \frac{a - c}{3}\cdot\frac{a - c}{3b} = \frac{(a - c)^2}{9b}.
\]
\textbf{Uniqueness.} The interior equations above are linear with a single solution. A profile with
$q_1 = 0$ would need $q_2 = BR_2(0) = k$ and then $BR_1(k) = k/2 > 0 = q_1$, a contradiction; the
same for $q_2 = 0$. Hence $(q^*, q^*)$ is the unique Nash equilibrium.

\answer{$q_1^* = q_2^* = \dfrac{a - c}{3b}$, with $Q^* = \dfrac{2(a-c)}{3b}$, $p^* = \dfrac{a + 2c}{3}$ and
profit $\dfrac{(a-c)^2}{9b}$ for each firm. The book's answer is the same.}

\interp Compare with the two benchmarks of the slides (Lecture~2, slide~22). A monopolist produces
$(a - c)/(2b)$ and earns $(a-c)^2/(4b)$; perfect competition has $p = c$, $Q = (a-c)/b$. The duopoly
sits in between: $Q^* = \tfrac23 (a-c)/b$ is more than the monopoly output $\tfrac12 (a-c)/b$ and less
than the competitive output, and industry profit $2(a-c)^2/(9b)$ is below the monopoly's
$(a-c)^2/(4b)$ ($0.222 < 0.25$ in units of $(a-c)^2/b$). Each firm ignores the harm its output
does to the other's price, so total output is too high from the cartel's point of view
(Lecture~2, slides~25--26).

\verify \texttt{check\_L2.py} differentiates symbolically, solves the two first-order conditions
with sympy and compares with the book's printed answer $q^* = (a-c)/(3b)$ (and its best-response
function $(a - c)/(2b) - q_j/2$). It also recomputes the case $a = 10$, $b = 1$, $c = 2$ (so
$q^* = 8/3$, profit $64/9$) by exact search over a grid of quantities, and confirms that the
best response is $0$ at $q_j = (a-c)/b$.

\mistakes
\begin{itemize}
  \item Maximising \emph{total} profit $\pi_1 + \pi_2$ instead of each firm's own profit: that is
  the cartel, $q_i = (a-c)/(4b)$.
  \item Forgetting that $q_j$ also enters the price, so the coefficient on $q_i$ in the first-order
  condition is $2b$, not $b$.
  \item Dropping the corner: the true best response is $\max\{0, \cdot\}$.
\end{itemize}
\end{exsol}
